% ============================================================================
\chapter{TCP/IP 协议栈}
\label{ch:tcpip}
% Stage H-4 ~ H-6
% ============================================================================

\section{本章目标}

\begin{itemize}
  \item 理解协议栈分层架构与数据封装/解封
  \item 设计可插拔 \texttt{proto\_ops\_t} 接口 + 协议注册表
  \item 设计 \texttt{sk\_buff} 网络缓冲区
  \item 实现 ARP + IPv4 + ICMP（能 ping 通 QEMU 网关）
  \item 实现 UDP + TCP（三次握手 / 数据传输 / 四次挥手）
\end{itemize}

% ============================================================================
\section{概念：协议栈分层}
\label{sec:tcpip-concept}
% ============================================================================

% TODO:
% - 自底向上数据流：
%     网卡收到以太网帧 → Link 层解析 EtherType
%     → Network 层解析 IP 头 → 路由/转发
%     → Transport 层解析 TCP/UDP 头 → 递交给 Socket
% - sk_buff（网络缓冲区）设计：
%     一块连续内存，通过 head/data/tail/end 指针在各层之间共享
%     push/pull 操作添加/剥离协议头
% - 协议注册表：根据 EtherType/IP Protocol/Port 分发到不同处理函数

% ============================================================================
\section{可插拔接口：\texttt{proto\_ops\_t}}
\label{sec:tcpip-interface}
% ============================================================================

% TODO:
% - 接口设计：
%     typedef struct proto_ops {
%         const char* name;                    // "tcp", "udp", "icmp"
%         int (*send)(struct sk_buff* skb);
%         int (*recv)(struct sk_buff* skb);
%         int (*bind)(uint16_t port);
%         int (*connect)(uint32_t ip, uint16_t port);
%     } proto_ops_t;
% - 协议注册：proto_register(IPPROTO_TCP, tcp_get_ops())
% - IP 层根据 protocol 字段分发到对应 proto_ops

% ============================================================================
\section{ARP + IPv4 + ICMP}
\label{sec:tcpip-arp-ip-icmp}
% ============================================================================

% TODO:
% - ARP (Address Resolution Protocol)：
%     IP 地址 → MAC 地址映射
%     ARP 请求：广播 "Who has 10.0.2.2?"
%     ARP 应答：单播 "10.0.2.2 is at xx:xx:xx:xx:xx:xx"
%     ARP 缓存表（IP → MAC，带超时）
% - IPv4：
%     IP 头结构：version(4), IHL, TTL, protocol, src_ip, dst_ip, checksum
%     路由表（简单：单默认网关）
%     checksum 计算（one's complement sum）
% - ICMP (Internet Control Message Protocol)：
%     Echo Request (type=8) / Echo Reply (type=0)
%     ping 实现：构造 echo request → 发送 → 等待 reply → 计算 RTT
% - 验证：ping 10.0.2.2（QEMU user-net 网关）

% ============================================================================
\section{UDP + TCP}
\label{sec:tcpip-tcp}
% ============================================================================

% TODO:
% - UDP (User Datagram Protocol)：
%     无连接、不可靠、简单
%     UDP 头：src_port, dst_port, length, checksum
%     sendto/recvfrom 语义
% - TCP (Transmission Control Protocol)：
%     面向连接、可靠传输、流控制
%     TCP 头：src_port, dst_port, seq, ack, flags(SYN/ACK/FIN/RST), window, checksum
%     三次握手：SYN → SYN+ACK → ACK
%     四次挥手：FIN → ACK → FIN → ACK
%     TCP 状态机图（CLOSED→LISTEN→SYN_RCVD→ESTABLISHED→...→TIME_WAIT）
%     滑动窗口（简化版：固定窗口大小）
%     超时重传（简单定时器，指数退避）
% - 验证：TCP 连接到 QEMU host 上的 echo server

% ============================================================================
\section{验证}
\label{sec:tcpip-verify}
% ============================================================================

% TODO:
% - 测试 1：ARP 请求 → 收到应答 → 缓存表正确
% - 测试 2：ping 10.0.2.2 → ICMP reply → RTT 显示
% - 测试 3：UDP sendto → host 收到数据
% - 测试 4：TCP 三次握手 → 发送数据 → 收到回复 → 四次挥手
% - shell 命令：ping / arp / net test

% ============================================================================
\section{本章小结}
% ============================================================================

TCP/IP 协议栈通过可插拔 \texttt{proto\_ops\_t} 接口，让协议实现可独立扩展。
ARP + IP + ICMP 构成基本网络通信能力，TCP + UDP 提供可靠和不可靠传输选择。
下一章实现 BSD Socket API，将网络能力暴露给用户态程序。
